% 对应题号 M1-26，复用已核对的本题几何图。
% 修订：启用 xeCJK 自动换行；保留原题参数与物理关系。
\documentclass[tikz,border=14pt]{standalone}
\usepackage{fontspec}
\usepackage{xeCJK}
\usepackage{amsmath,amssymb}
\usepackage{xcolor}
\usepackage{tikz}
\usetikzlibrary{arrows.meta,calc,patterns,positioning,bending}
\IfFontExistsTF{Heiti SC}{\IfFontExistsTF{Heiti SC}{\setmainfont{Heiti SC}}{\setmainfont{FandolHei-Regular}}\IfFontExistsTF{Heiti SC}{\setsansfont{Heiti SC}}{\setsansfont{TeX Gyre Heros}}}{\IfFontExistsTF{Songti SC}{\setmainfont{Songti SC}}{\setmainfont{FandolSong-Regular}}\IfFontExistsTF{Songti SC}{\setsansfont{Songti SC}}{\setsansfont{FandolSong-Regular}}}
\definecolor{bluefield}{RGB}{24,76,255}
\definecolor{redfield}{RGB}{232,49,35}
\definecolor{orangefield}{RGB}{236,132,18}
\definecolor{greenfield}{RGB}{18,132,86}
\definecolor{purplefield}{RGB}{112,72,190}
\definecolor{linegray}{RGB}{60,68,78}
\definecolor{softblue}{RGB}{238,243,255}
\definecolor{softgreen}{RGB}{238,249,244}
\definecolor{softorange}{RGB}{255,247,235}
\tikzset{
  arrow/.style={-{Latex[length=3.1mm,width=2.1mm]},line width=1.1pt},
  thinarr/.style={-{Latex[length=2.5mm,width=1.7mm]},line width=.85pt},
  axis/.style={-{Latex[length=2.5mm,width=1.7mm]},line width=.8pt,draw=linegray},
  guide/.style={draw=linegray!55,dashed,line width=.65pt},
  box/.style={draw=linegray,rounded corners=3pt,fill=white,align=center,inner sep=6pt,line width=.75pt},
  bluebox/.style={box,draw=bluefield,fill=softblue},
  greenbox/.style={box,draw=greenfield,fill=softgreen},
  orangebox/.style={box,draw=orangefield,fill=softorange}
}
\IfFontExistsTF{Heiti SC}{\setCJKmainfont[AutoFakeBold=1.5]{Heiti SC}}{\setCJKmainfont[AutoFakeBold=1.5]{FandolHei-Regular}}
\IfFontExistsTF{Heiti SC}{\setCJKsansfont{Heiti SC}}{\setCJKsansfont{FandolHei-Regular}}
\begin{document}

\begin{tikzpicture}[x=1cm,y=1cm]
\path[use as bounding box] (-10,-4.6) rectangle (10,4.6);
\fill[white] (-10,-4.6) rectangle (10,4.6);
\node[font=\Large\bfseries,text width=18.6cm,align=center] at (0,4.05) {汤川型势的保守力：方向径向向外，近区类库仑，远区指数屏蔽};
\begin{scope}[shift={(-5.0,-.15)}]
\foreach \r in {.75,1.45,2.15}{\draw[linegray!45] (0,0) circle(\r);}
\fill[redfield] (0,0) circle(.12) node[below=5pt,fill=white,inner sep=1pt] {相互作用中心};
\fill[bluefield] (2.0,1.1) circle(.1) node[above left=7pt,fill=white,inner sep=1pt] {试探粒子};
\draw[arrow,linegray] (0,0)--(2.0,1.1) node[midway,above] {$\mathbf r$};
\draw[arrow,greenfield] (2.0,1.1)--(3.2,1.75) node[right] {$\mathbf F$};

\node[orangebox,text width=5.8cm] at (0,-2.8)
  {$U_0>0$ 且 $dU/dr<0$，故 $-\nabla U$ 沿 $+\hat{\mathbf r}$，表现为排斥。};
\end{scope}
\begin{scope}[shift={(4.7,-.15)}]
\draw[axis] (-3,0)--(3.2,0) node[right] {$r$};
\draw[axis] (-2.7,-.2)--(-2.7,2.8) node[above] {$|\mathbf F|$};
\draw[redfield,line width=1.2pt] plot[domain=-2.25:2.9,samples=100]
  (\x,{1.7*exp(-.65*(\x+2.25))/(.75+(\x+2.25)^2)});
\node[bluebox,text width=6.5cm] at (0,-1.3)
  {$\displaystyle \mathbf F=U_0e^{-\alpha r}\left(\frac1{r^2}+\frac{\alpha}{r}\right)\hat{\mathbf r}$};
\node[greenbox,text width=6.5cm] at (0,-2.8)
  {$r\ll\alpha^{-1}$：$|\mathbf F|\sim U_0/r^2$；\\
   $r\gg\alpha^{-1}$：$|\mathbf F|\sim\alpha U_0e^{-\alpha r}/r$。};
\end{scope}
\end{tikzpicture}
\end{document}
